\mode*

\subsection{Elimination, Relevanter Konzentrationsabfall}


\frame{
\mode<presentation>{
\begin{tikzpicture}[remember picture, overlay]
\node[inner sep=0pt,xshift=0cm] at (current page.center){
\tikz \draw[step=2mm,black!50] (0,0) grid (126mm,94mm);
};
{
\node[coordinate] at (1cm,-3.55cm) (start) {};
\node[coordinate] at (1cm,3cm) (end) {};
\node[coordinate] at (10cm,-3.55cm) (endy) {};
\draw[->,thick] (start)--(end);
\draw[->,thick] (start)--(endy);
};
\end{tikzpicture}
}
\frametitle{Klinisch relevanter Konzentrationsabfall}
\note<1>{
\begin{itemize}
\item
Sigmoide Kurve: Keine fixe Konzentration, keine fixe Zeit: Wenn Konzentration um 70\% abfallen muss damit Patient z.B. erwacht, dann ist diese Zeit relevant.
\item
Excel Simulation - Vergleich von verschiedenen Kurven.
\end{itemize}
}
\infina
}


\mode<article>{
\vspace{1cm}
\begin{tikzpicture}[auto]
\node[inner sep=0pt,xshift=0cm] at (current page.center){
\tikz \draw[step=2mm,black!50] (0,0) grid (140mm,70mm);
};
\end{tikzpicture}

\vspace{0.5cm}
Um abzusch�tzen, wie lange es nach abstellen der Medikamentenzufuhr geht, bis die Wirkung \enquote{weg ist} muss man wissen, wieviel die Konzentration absinken muss. Es gibt keinen Grund anzunehmen, dass f�r die verschiedenen Wirkungen der An�sthetika, ein Abfall der Konzentration auf die H�lfte (\HWZ !) klinisch klinisch relevant ist. F�r volatile An�sthetika braucht es Reduktion der Konzentration um ca. 70 - 80 \% damit ein Patient erwacht. (ad�quate An�sthesie bis Erwachen). Eine �hnliche Gr�ssenordnung ist auch f�r Propofol beschrieben. 



}